\source{Лекции, с. 50–51; Карпов, §14.1; пропущенные шаги индукции раскрыты составителем}
\begin{definition}
Пусть \(G\) простой, \(n\) — число вершин, \(m\) — число рёбер,
\(c=c(G)\). Запишем
\[
 \chi_G(k)=\sum_{j=0}^n(-1)^j b_j(G)k^{n-j}.
\]
Знакопеременность означает \(b_j\ge0\); ниже уточняется, какие из них строго положительны.
\end{definition}
\begin{theorem}[Старшие коэффициенты]
Старший коэффициент \(\chi_G\) равен 1, а при \(n\ge1\) коэффициент
при \(k^{n-1}\) равен \(-m\).
\end{theorem}
\begin{proof}
Индукция по числу рёбер, одновременно для всех чисел вершин.
При \(m=0\) имеем \(\chi_G(k)=k^n\); второй коэффициент нулевой.
Пусть \(e\in E\). По индукции
\[
\chi_{G-e}(k)=k^n-(m-1)k^{n-1}+\text{члены меньшей степени},
\]
а \(\chi_{G/e}\) унитарен степени \(n-1\) по билету 18.
Из \(\chi_G=\chi_{G-e}-\chi_{G/e}\) получаем старший коэффициент 1
и следующий \(-(m-1)-1=-m\). Кратные рёбра после стягивания
убираются, что не влияет на унитарность; сравнивать их число здесь не требуется.
При \(n=1\) имеем \(m=0\) и \(\chi_G=k\).
\end{proof}
\begin{lemma}[Рекурсия для модулей коэффициентов]
Для ребра \(e\)
\[
 b_j(G)=b_j(G-e)+b_{j-1}(G/e),
\]
где коэффициенты с индексами вне допустимого диапазона считаются нулевыми,
а кратные рёбра после стягивания удаляются.
\end{lemma}
\begin{proof}
Подставляем разложения многочленов в
\(\chi_G=\chi_{G-e}-\chi_{G/e}\). У второго графа на одну вершину меньше,
поэтому его коэффициент при \(k^{n-j}\) равен
\((-1)^{j-1}b_{j-1}(G/e)\). Вычитание меняет знак на \((-1)^j\),
и после сравнения коэффициентов получается сумма.
\end{proof}
\begin{theorem}[Чередование знаков]
\[
 b_j(G)>0\quad(0\le j\le n-c),\qquad
 b_j(G)=0\quad(j>n-c).
\]
Таким образом, ненулевые коэффициенты при степенях
\(k^n,k^{n-1},\ldots,k^c\) строго чередуются по знаку,
а коэффициенты более низких степеней нулевые.
\end{theorem}
\begin{proof}
Индукция по \(n+m\). Для графа без рёбер \(\chi_G=k^n\), \(c=n\),
так что положителен только \(b_0=1\).
Пусть есть ребро \(e\), положим \(r=n-c\ge1\).
Стягивание не меняет число компонент, следовательно,
\((n-1)-c(G/e)=r-1\). При удалении немоста число компонент не меняется,
а при удалении моста увеличивается на 1.
По индукции оба слагаемых рекурсии неотрицательны.
Если \(e\) не мост, первое слагаемое положительно для \(0\le j\le r\)
и равно нулю выше \(r\); второе может быть положительно только при
\(1\le j\le r\). Если \(e\) мост, первое положительно при
\(0\le j\le r-1\), второе при \(1\le j\le r\).
В обоих случаях сумма положительна ровно для \(0\le j\le r\).
\end{proof}
\begin{example}
Для дерева на четырёх вершинах
\(\chi_G(k)=k(k-1)^3=k^4-3k^3+3k^2-k\).
Старший коэффициент 1, второй \(-3=-m\), ненулевые знаки чередуются;
свободный член равен нулю.
\end{example}

\begin{cor}[Наименьшая ненулевая степень]
Наименьшая степень с ненулевым коэффициентом в \(\chi_G(k)\) равна \(c(G)\).
В частности, непустой граф связен тогда и только тогда,
когда коэффициент при \(k\) ненулевой.
\end{cor}
\begin{proof}
В теореме о чередовании последний положительный модуль — \(b_{n-c}\),
стоящий при \(k^c\); все следующие коэффициенты нулевые.
\end{proof}
\textbf{Дополнительное доказательство по включениям–исключениям.}
Следующая формула отсутствует в скане и не нужна для лекционной индукции.
\begin{lemma}[Включения–исключения]
\[
 \chi_G(k)=\sum_{S\subseteq E}(-1)^{|S|}k^{c(V,S)},
\]
где \(c(V,S)\) — число компонент остовного подграфа с рёбрами \(S\).
\end{lemma}
\begin{proof}
Из всех \(k^n\) раскрасок исключаем события \(B_e\): концы ребра \(e\)
имеют одинаковый цвет. При выполнении всех событий для \(e\in S\)
вершины каждой компоненты \((V,S)\) должны быть одноцветны; цвета
разных компонент выбираются независимо, поэтому пересечение событий
имеет мощность \(k^{c(V,S)}\). Формула включений–исключений даёт нужную сумму.
Для полноты: раскраска с \(t\) нарушенными рёбрами получает в этой сумме
вес \(\sum_{r=0}^t(-1)^r\binom tr=(1-1)^t\), равный 1 при \(t=0\)
и 0 при \(t>0\), что непосредственно доказывает формулу.
\end{proof}
